% UTM license thesis/project starter.
% Replace every red placeholder with verified details from the approved task sheet,
% your own research and the supervisor's instructions before preparing a final copy.
% Check the course curriculum for clear permission before using generative AI.
% If permitted, disclose its use in the form required by the instructor or department.
% For a thesis written in another language, the UTM guide requires title pages in
% that language and Romanian. Set the institutional fields and language before each
% \makeutmthesistitlepage call; the class counts both title pages automatically.
\documentclass[thesis]{utmreport}
\selectlanguage{romanian}

\utmlogo{utm-logo.png}
\institution{Universitatea Tehnică a Moldovei}
\faculty{Facultatea Calculatoare, Informatică și Microelectronică}
\department{\placeholder{Departamentul de profil}}
\documenttype{Proiect / teză de licență}
\reporttitle{\placeholder{Titlul aprobat al proiectului / tezei de licență}}
\student{\placeholder{Prenumele NUMELE}}{\placeholder{Grupa academică}}
\supervisor{\placeholder{Prenumele NUMELE}}{\placeholder{titlul științifico-didactic și științific}}
\consultant{\placeholder{Prenumele NUMELE}}{\placeholder{titlul consultantului}}
\cityyear{Chișinău}{\placeholder{anul susținerii}}

\begin{document}
\makeutmcover
\makeutmthesistitlepage
\utmfrontmatter

\unnumberedchapter{Caiet de sarcini}
\placeholder{Introduceți aici paginile caietului de sarcini aprobat și completat de departament. Nu inventați tema, datele inițiale, termenele sau semnăturile.}
\clearpage

\unnumberedchapter{Declarația studentului}
\placeholder{Înlocuiți această pagină cu declarația oficială scrisă de mână în exemplarul final. Folosiți textul și instrucțiunile cerute de UTM și de departament.}
\clearpage

\unnumberedchapter{Avizul coordonatorului}
\placeholder{Introduceți avizul completat și semnat de coordonator. Nu completați și nu semnați acest document în numele coordonatorului.}
\clearpage

\begin{utmabstract}[Rezumat]
\placeholder{Scrieți un rezumat de 300-500 de cuvinte despre scopul real, obiectivele, metodele și rezultatele obținute. Includeți numele autorului, titlul aprobat și structura lucrării.}

\utmkeywords{\placeholder{până la 5 cuvinte-cheie, cu litere mici}}
\end{utmabstract}

\selectlanguage{english}
\begin{utmabstract}[Abstract]
\placeholder{Provide an English translation of the Romanian abstract, including the title, with the same 300-500 word scope and only verified research details.}

\utmkeywords{\placeholder{up to 5 lowercase keywords}}
\end{utmabstract}
\selectlanguage{romanian}
% If the thesis language is neither Romanian nor English, add its original-language
% annotation too, following the current department and supervisor instructions.

\tableofcontents

% Include this list if the thesis uses abbreviations or terms outside Romanian.
\begin{utmabbreviations}
\item[\placeholder{API}] \placeholder{Interfață de programare a aplicațiilor}
\end{utmabbreviations}

\utmintroduction
\placeholder{Prezentați tema aprobată, actualitatea ei, problema, scopul, obiectivele, metodologia și structura lucrării, pe baza cercetării efectuate.}

\chapter{Cadrul teoretic}
\placeholder{Sintetizați critic sursele relevante pentru tema aprobată. Citați fiecare sursă în ordinea primei apariții, de exemplu \cite{source-a}.}

\chapter{Metodologia și rezultatele cercetării}
\placeholder{Descrieți metodele folosite, soluția realizată și rezultatele verificate. Explicați limitările și nu prezentați exemple ca măsurători.}

\utmconclusions
\placeholder{Formulați concluzii legate de obiective și rezultate, apoi indicați limitele și direcțiile următoare.}

\begin{thebibliography}{99}
\bibitem{source-a} \placeholder{Completați referința în stilul IEEE cerut de UTM. Citați-o în text ca [1].}
\end{thebibliography}

\appendix
\chapter{Materiale suplimentare}
\placeholder{Includeți aici numai materiale suplimentare la care faceți trimitere în text.}

\end{document}
